% Extended paper (journal). Builds on the JAIIO/ASAID 2026 camera-ready (the v1 study).
% Canonical model: M8 (hierarchical ordered-logit + model x instruction interaction).
\documentclass{llncs}

\usepackage[T1]{fontenc}
\usepackage{graphicx}
\usepackage{amsmath}
\usepackage{amssymb}
\usepackage{booktabs}
\usepackage{microtype}
\usepackage[htt]{hyphenat}
\emergencystretch=1.5em
\graphicspath{{figures/}{../report-figures/v6/}}
\usepackage[style=apa]{biblatex}
\addbibresource{references.bib}

\begin{document}

\title{Not All Instructions Are Forgotten Equal: Per-Instruction
Retention Is a Model-by-Instruction Interaction}
\titlerunning{Not All Instructions Are Forgotten Equal}
\author{Tom\'as P. Korenblit\orcidID{0009-0002-5682-8475}}
\authorrunning{T. P. Korenblit}
\institute{Universidad Nacional de San Mart\'in, Buenos Aires, Argentina\\
\email{tomaskorenblit@gmail.com}}
\maketitle

\begin{abstract}
Large language models lose adherence to user-stated preferences over long
conversations, and the common defense, restating every instruction each
turn, is expensive and indiscriminate. Whether a selective alternative is
possible depends on how retention is structured across instructions and
models. We measure per-instruction compliance in realistic multi-turn
coding sessions across five models, twelve instruction types, and three
open-source Python codebases, graded by deterministic checkers, and we
compare a ladder of Bayesian ordered-logistic models. The decisive
structure is neither a per-instruction main effect nor a per-model main
effect but their \emph{interaction}: a model that adds a
model-by-instruction interaction term improves out-of-sample fit by about
59 ELPD (five standard errors) over additive models, while additive
per-model terms add nothing. Which instruction is retained depends on the
model. We also find that the ``backfire'' reported in the conference version
of this study, where stating one instruction lowered compliance, does not
survive: it is a measurement artifact, confirmed both by a corrected checker
and by a negation-versus-positive framing experiment that finds no effect of
polarity. We argue that the recommended model, a hierarchical Bayesian
graded-response model with a model-by-instruction interaction, is the
appropriate analysis for evaluation data generally, and that reporting raw
means or additive effects mischaracterizes what evaluations measure. The
finding matters most for agents deployed in long-running loops without human
oversight, where a guardrail that decays does so silently.
\keywords{instruction retention \and item response theory \and Bayesian
inference \and model evaluation \and AI safety}
\end{abstract}

\section{Introduction}
\label{sec:intro}

Language models are increasingly deployed as autonomous agents that run in
loops, call tools, and act over hundreds of turns with no human reviewing
each step. An agent is governed by its instructions: the operational limits
and safety guardrails it is given, usually once and at the start, that must
hold for the whole run. Yet adherence to instructions erodes as the context
grows, a phenomenon we refer to as \emph{context rot}. \textcite{laban2025}
measured average compliance drops of 39\% across 15 models in multi-turn
settings. The cause is architectural: attention is not uniformly
distributed across the context window, and information in middle positions
receives less weight than at the boundaries \parencite{liu2024,
chowdhury2026}. An instruction stated once and then buried under
accumulating tool output competes for attention against everything that
follows it, and in a deployed loop there is no user to notice the lapse and
restate it.

This is a safety concern before it is a cost concern. When the model drops a
preference, a chat user notices and restates it; a deployed agent has no
such supervisor, so a guardrail that decays does so silently while the agent
keeps acting through tools whose effects are real and often irreversible
\parencite{mu2025}. Repeating every instruction each turn, the standard
mitigation, multiplies prompt cost and still cannot tell an operator which
constraint remains in force. Because safety is conjunctive, an aggregate
compliance rate says nothing about whether the one rule that matters is
among those that lapsed.

Whether reinforcement can be targeted instead of uniform depends on how
retention is structured. The conference version of this work
\parencite{korenblit2026} fit a per-instruction Bayesian model to a
single-model dataset and reported that retention varies widely across
instructions, that one instruction was harmed by being stated, and that the
model and codebase explained little variance. This paper revisits those
claims with a larger, multi-model dataset, a corrected scoring pipeline, a
framing experiment, and a wider model comparison. The picture changes.

Our contributions are: (i) we show that per-instruction retention is
governed by a \emph{model-by-instruction interaction}, not by a main effect
of either factor: an interaction model improves leave-one-out fit by about
59 ELPD over additive models, and additive per-model terms add nothing;
(ii) we show that the reported ``backfire'' is a measurement artifact, with
two independent lines of evidence (a corrected checker and a
polarity-framing experiment); (iii) we quantify how fragile
variance-component conclusions are to the pool of models analyzed, and which
conclusions survive; and (iv) we argue, and demonstrate on this dataset,
that a hierarchical Bayesian graded-response model with a
model-by-instruction interaction is the appropriate default for analyzing
evaluation data, situating it against item-response-theory and
generalizability-theory practice.

\section{Related work}
\label{sec:related}

\paragraph{Instruction degradation in long contexts.}
Several studies establish that adherence falls as conversations grow.
\textcite{laban2025} report an average 39\% compliance drop across 15
models; \textcite{he2024} show multi-instruction accuracy falling within
three turns; \textcite{du2025} find that a larger context window does not by
itself fix the loss. None asks whether different instructions are retained
to different degrees. On mechanism, \textcite{liu2024} document a U-shaped
positional attention curve and \textcite{mu2025} a saturation effect.
Existing mitigations reinforce uniformly: periodic reminders
\parencite{dongre2025} and prompt repetition \parencite{leviathan2025},
which can be counterproductive for instructions where restating degrades
compliance.

\paragraph{Evaluation as measurement, and item response theory.}
A parallel literature argues that benchmark scores should be treated as
measurements with uncertainty rather than raw means.
\textcite{miller2024errorbars} imports experiment-analysis machinery
(standard errors, clustering, paired differences); \textcite{bowyer2025clt}
shows the normal approximation fails at the small sample sizes typical of
specialized evals. Construct-validity work \parencite{bean2025construct}
finds widespread operationalization failures across benchmarks, and
generalizability theory has been applied to decompose LLM-scoring variance
across facets \parencite{song2025gtheory, brennan2001gtheory}. Ranking is
increasingly treated as estimation, with Bradley--Terry models and
confidence intervals \parencite{chiang2024arena}. Item response theory
casts systems as examinees with a latent ability and items as having
difficulty and discrimination \parencite{lalor2016irt,
deboeck2004explanatory, samejima1969grm}; it has been used to build
evaluation scales, reweight leaderboards
\parencite{rodriguez2021leaderboards}, compare test sets
\parencite{vania2021testsets}, and build efficient benchmarks
\parencite{polo2024tinybenchmarks}. This work is almost entirely on binary
correct/incorrect outcomes and estimated by variational approximation; the
one graded-response application we are aware of targets judge reliability,
not the systems under test \parencite{choi2026judgeirt}. Our model sits in
the explanatory-IRT tradition \parencite{deboeck2004explanatory} with
crossed random effects \parencite{baayen2008crossed}, but treats responses
as ordinal, casts the experimental manipulation as a covariate on ability,
and estimates the full hierarchy by MCMC.

\paragraph{Tracking what a model has forgotten.}
Bayesian Knowledge Tracing \parencite{corbett1994} estimates per-concept
mastery from response sequences. We adapt the framing: the model is the
student and its instructions are the concepts. A reinforcement scheme would
only be worthwhile if retention is heterogeneous and predictable, which is
what we test.

\section{Method}
\label{sec:method}

\subsection{Decision planting}
We simulate a coding assistant working on a real open-source codebase over
25 turns. Real source files are injected at turns 0--19, building context
pressure. In the \textbf{treatment} condition, 12 coding preferences are
embedded as casual inline requests in the user's messages, six early
(turns 1--7) and six late (turns 14--19), spanning code style, architecture,
testing, naming, dependencies, and documentation. They are phrased as
natural user preferences, not system-prompt rules (e.g., ``Important: when
writing array operations for this project, always use NumPy broadcasting
instead of for loops''). The \textbf{baseline} condition is identical but
plants nothing, measuring spontaneous compliance against the same probes. At
turns 20--25 the user requests code-generation tasks that elicit each
decision; each test turn probes two decisions, one early and one late.

To separate an instruction's content from its phrasing, two further arms
hold meaning constant and vary only polarity: \textbf{positive} and
\textbf{negation} paraphrases of each flippable instruction. The framing
effect is estimated arm against arm, since the canonical treatment wording
differs from both by more than polarity.

\subsection{Compliance measurement}
Each decision is scored by a deterministic checker that parses the generated
code with pattern matching, with no LLM in the scoring loop, so results are
reproducible across runs. Each checker emits an ordinal 0--3 score
(0 = ignored, 3 = fully followed) or abstains ($-1$, excluded) when the code
gives no opportunity to apply the rule. The conference study's
\texttt{no\_new\_imports} checker counted any import statement as new; the
corrected checker treats re-showing an already-imported module as compliant
(\S\ref{sec:backfire}). As a secondary channel, a panel of three
cross-family LLM judges scores the same outputs blind to condition and
model; no judge shares a family with a candidate, removing the conference
study's self-judging bias. We report checker--judge agreement in
\S\ref{sec:validity}.

\subsection{A ladder of Bayesian models}
\label{sec:models}
We model the ordinal scores with hierarchical ordered-logistic
(cumulative-link) regression, comparing five specifications of increasing
structure by leave-one-out cross-validation \parencite{vehtari2017loo}. Let
$d$ index decision, $m$ model, $k$ codebase, and let $t_i$ be the treatment
indicator. All effects use non-centered parameterizations and offset
cutpoints; priors follow standard weakly-informative choices for ordinal
models \parencite{mcelreath2020rethinking, burkner2021irt}.

\textbf{M0} (the conference model) is per-decision only:
$\eta_i = \alpha_{d} + \beta_{d}\, t_i$, with $\alpha_d, \beta_d$ given
hierarchical priors. \textbf{M4} adds additive per-model and per-codebase
intercepts and a per-model treatment slope:
$\eta_i = \alpha_d + \beta_d t_i + g_m + t_i\, s_m + g_k$. \textbf{M\_corr}
draws $(g_m, s_m)$ from a correlated bivariate normal (LKJ prior).
\textbf{M\_irt} is an explanatory graded-response model: it replaces the
additive per-model intercept with a multiplicative discrimination term,
$\eta_i = b_d + a_d\,\theta_m + \beta_d t_i + g_k$, where $\theta_m$ is a
latent model ability and $a_d>0$ a per-instruction discrimination, the IRT
form of a model-by-instruction interaction. \textbf{M8} keeps the additive
structure of M4 and adds an identifiable interaction \emph{variance
component}: a per-cell treatment effect
$z_{m,d}\sim\mathcal{N}(0, \sigma_{\text{int}})$, so
$\eta_i = \alpha_d + \beta_d t_i + g_m + t_i s_m + g_k + t_i\, z_{m,d}$. The
key scientific parameter shifts from $\sigma_\beta$ (spread of
per-instruction effects) to $\sigma_{\text{int}}$ (spread of the
model-by-instruction interaction). A variance-component interaction is
preferred to a rank-1 factor model, which is unidentified up to sign and
scale; the factor version did not converge ($\hat R = 1.74$), while M8
samples cleanly. Fit in PyMC \parencite{pymc2023} with NUTS via nutpie
\parencite{nutpie2024, hoffman2014}, four chains, 2000 draws each,
\texttt{target\_accept} $= 0.99$.

\section{Experimental setup}
\label{sec:setup}
Three open-source Python libraries span a range of context pressure: Bambi
(median peak $\approx$85K tokens by the test turns), ArviZ ($\approx$149K),
and PyMC ($\approx$187K). Five instruction-following models were evaluated
with both baseline and treatment in every codebase: Qwen 3.5 (27B),
Gemini 3.1 Flash Lite, GPT-5.4-nano, Nemotron 3 Super (120B), and
Gemma 4 (26B). Each condition was run for five epochs to average per-call
stochasticity. After excluding abstentions, the primary dataset is 1{,}067
compliance observations across 5 models, 12 decision types, 3 codebases, and
2 conditions; the framing arms add a further set across three of the models.
The pipeline is a port to the Inspect evaluation framework; the conference
study used a bespoke runner with a single baseline model and a same-family
judge.

\section{Results}
\label{sec:results}

\subsection{Retention is heterogeneous}
The per-instruction treatment effect varies substantially: in the additive
model its group-level standard deviation is $\sigma_\beta = 1.19$ (94\% HDI
$[0.55, 1.99]$), excluding zero. Five to seven of the twelve instructions,
depending on specification, have a treatment effect whose interval excludes
zero (Fig.~\ref{fig:beta}); the strongest are
\texttt{dependencies\_module\_constants}, \texttt{docs\_regex\_comments},
\texttt{testing\_parametrize}, and \texttt{docs\_numpy\_style}. This is
about half the magnitude reported by the conference study
($\sigma_\beta = 2.11$); \S\ref{sec:interaction} explains why.

\begin{figure}[t]
\centering
\includegraphics[width=0.78\textwidth]{fig_beta_forest}
\caption{Per-instruction treatment effect $\beta$ with 94\% HDI (additive
model, clean five-model set). Gold: interval above zero; grey: crosses
zero. The conference study's $-2.36$ ``backfire'' on
\texttt{dependencies\_no\_new} is gone (bottom).}
\label{fig:beta}
\end{figure}

\subsection{The structure is a model-by-instruction interaction}
\label{sec:interaction}
Comparing the model ladder by leave-one-out cross-validation
(Table~\ref{tab:loo}, Fig.~\ref{fig:loo}) gives a single clear result. The
two models that capture a model-by-instruction interaction, M8 and M\_irt,
fit decisively better than every additive model, by about 59 ELPD (more
than five standard errors). The two interaction models are not
distinguishable from each other ($\Delta$ELPD $= 9.5 \pm 10.5$), so it is
the interaction that matters, not its exact parameterization. Crucially, the
additive per-model terms add nothing: M4 (per-model intercept and slope) is
indistinguishable from M0 (per-decision only), and allowing the model
intercept and slope to correlate (M\_corr) does not help; the estimated
correlation is $-0.37$ (94\% HDI $[-0.97, 0.35]$), so a model that complies
more by default is not more reinforceable.

\begin{table}[t]
\centering
\caption{Model comparison by leave-one-out cross-validation (clean
five-model set, 1{,}067 observations). Higher ELPD is better;
$\Delta$ELPD is the difference from the best model with its standard
error.}
\label{tab:loo}
\begin{tabular}{llrr}
\toprule
Model & Interaction structure & ELPD & $\Delta$ELPD (SE) \\
\midrule
M8 & model$\times$instruction variance & $-545.9$ & --- \\
M\_irt & 2PL discrimination $a_d\theta_m$ & $-555.4$ & $9.5\ (10.5)$ \\
M\_corr & correlated per-model slopes & $-604.4$ & $58.5\ (11.2)$ \\
M4 & additive per-model & $-604.8$ & $58.9\ (11.2)$ \\
M0 & per-instruction only & $-607.2$ & $61.3\ (11.7)$ \\
\bottomrule
\end{tabular}
\end{table}

\begin{figure}[t]
\centering
\includegraphics[width=0.85\textwidth]{fig_ladder_loo}
\caption{Out-of-sample fit (ELPD, higher is better) with standard errors.
The interaction models (teal) separate cleanly from the additive models
(grey), which cluster together.}
\label{fig:loo}
\end{figure}

Under M8 the interaction dominates: $\sigma_{\text{int}} = 1.84$ (94\% HDI
$[1.26, 2.40]$), larger than the per-instruction spread
($\sigma_\beta = 0.43$), the per-model intercept ($\sigma_{\text{model}} =
0.66$), and the per-model slope ($\sigma_{t\text{-model}} = 0.43$); see
Fig.~\ref{fig:m8var}. The additive components shrink once the interaction is
in the model, because in additive specifications they absorb interaction
variance. This also explains the halved $\sigma_\beta$ relative to the
conference study: with a single model, the model-by-instruction interaction
has nowhere to go and collapses into the per-instruction effect, inflating
it. Fig.~\ref{fig:heat} shows the interaction directly: the conference
study's primary model, Qwen, carries the most extreme cells, including the
strongly negative \texttt{dependencies\_no\_new} cell that produced the
original ``backfire.''

\begin{figure}[t]
\centering
\includegraphics[width=0.7\textwidth]{fig_m8_variance}
\caption{Group-level standard deviations under M8. The
model-by-instruction interaction $\sigma_{\text{int}}$ is the largest source
of structured variation; the additive per-model and per-codebase terms are
small.}
\label{fig:m8var}
\end{figure}

\begin{figure}[t]
\centering
\includegraphics[width=\textwidth]{fig_interaction_heatmap}
\caption{Posterior mean of the model-by-instruction interaction $z_{m,d}$
(log-odds). Red raises, blue lowers compliance for that pair beyond the
additive effects. Retention is a (model, instruction) property: which
instruction is retained depends on the model.}
\label{fig:heat}
\end{figure}

\subsection{The reported backfire is a measurement artifact}
\label{sec:backfire}
The conference study reported that stating \texttt{dependencies\_no\_new}
(``do not add new imports'') \emph{lowered} compliance, $\beta = -2.36$. Two
independent analyses show this is not real. First, the original checker
counted any import statement as new, including imports already present in the
files shown to the model; being told ``do not add imports'' plausibly makes
a model re-show existing imports, which the checker then penalized. With a
context-aware checker that only counts genuinely new modules, the effect is
null ($\beta = -0.11$ to $-0.25$ across specifications, interval spanning
zero). Second, the framing experiment finds no effect of polarity: the
pooled negation-minus-positive effect is $0.36$ (94\% HDI $[-0.13, 0.89]$),
and every per-instruction interval crosses zero, so phrasing the rule as a
negation does not reduce compliance. The artifact is instructive for
measurement: under M\_irt, \texttt{dependencies\_no\_new} has by far the
highest discrimination ($a_d = 2.90$), i.e.\ it separates models most
sharply, a reminder that a highly discriminating item can be measuring an
artifact rather than the intended construct.

\subsection{Variance components are pool-sensitive; rankings are not}
\label{sec:sensitivity}
The variance-component estimates depend on which models are pooled. Across
the five-model, six-model, and ten-model sets, $\sigma_\beta$ ranges
$1.0$--$1.7$ and $\sigma_{\text{model}}$ ranges $0.48$--$0.98$; the
interaction conclusion and the per-instruction \emph{rankings} are stable,
but the absolute magnitudes of the variance components are not. A
power-scaling sensitivity analysis confirms that the variance components are
prior-influenced (a consequence of having few groups), while the fixed
effects are not. We therefore report rankings and the interaction structure
as robust, and treat the variance magnitudes as estimates with genuine prior
dependence.

\subsection{Measurement validity}
\label{sec:validity}
Checker--judge agreement is weak: the quadratic-weighted $\kappa$ between the
deterministic checker and the cross-family judge panel has a panel median of
$0.13$, with individual judges from $0.03$ to $0.17$. Agreement is lowest on
judgment-laden style and documentation decisions and higher on syntactically
unambiguous ones, on which the robust effects sit. We read this as a limit on
construct validity for the softer decisions, not as support for replacing the
reproducible checker with a judge.

\section{Discussion}

The central result reframes the phenomenon. Per-instruction retention is not
a property of the instruction alone, nor of the model alone, but of the
(model, instruction) pair. An additive analysis, including the conference
model and any report of per-instruction means, misattributes interaction
variance to main effects and can manufacture artifacts like the
``backfire,'' which was a single extreme cell of one model amplified by a
checker bug. A reinforcement policy that follows from this is necessarily
model-specific: there is no instruction-level ranking that transfers across
models, because the ranking is itself a function of the model.

For deployed agents the implication is direct. A guardrail is an instruction,
and whether it survives a long run depends on the specific (model, guardrail)
pair, not on a generic decay rate. An operator cannot read off from one model
which constraints another model will drop, and an aggregate compliance number
hides the interaction entirely. Per-instruction, per-model monitoring is the
prerequisite for noticing silent decay.

The methodological lesson generalizes beyond retention. Evaluation data is
typically systems crossed with items, scored on graded rubrics, and the
quantity of interest is often comparative. The model that fits this data
well, a hierarchical Bayesian graded-response model with crossed effects and
a model-by-instruction interaction, is exactly what the item-response and
generalizability-theory traditions prescribe, yet current eval practice
reports raw means or, at best, additive uncertainty. Our ladder shows the
cost concretely: the additive models are 59 ELPD worse and miss the dominant
structure. We suggest this model as a default for eval analysis and develop
the argument in a companion paper.

\section{Limitations}
\label{sec:limitations}
The variance components are prior-sensitive and pool-dependent
(\S\ref{sec:sensitivity}); we report rankings and structure rather than
leaning on the magnitudes. The framing arms cover three of the five models,
so the polarity null is not full-factorial. Compliance is measured at turns
20--25 only, a retention snapshot, not a decay curve; we do not estimate
per-instruction forgetting rates. Checker--judge agreement is low on softer
decisions (\S\ref{sec:validity}), a construct-validity limit. The selective
reinforcement we motivate is not implemented as a closed-loop policy; we
estimate the structure such a policy would exploit, not its closed-loop
performance. Context length varies across codebases rather than being held
constant, though the per-codebase variance is small.

\section{Conclusion}
Per-instruction retention in long coding conversations is governed by a
model-by-instruction interaction: an interaction model fits about 59 ELPD
better than additive models, additive per-model terms add nothing, and which
instruction is retained depends on the model. The conference study's
single-model design collapsed this interaction into an inflated
per-instruction effect and, together with a checker artifact, produced a
spurious ``backfire'' that does not replicate. The heterogeneity is real but
smaller and differently shaped than first reported. Beyond the retention
question, the analysis argues for a specific standard: evaluation data with
systems crossed against graded items should be analyzed with a hierarchical
Bayesian interaction model, not raw means.

\subsubsection*{Use of large language models.}
During the preparation of this work the author used a large language model
(Claude) to improve the language and readability of the manuscript and to
assist with analysis and figure code; all results were computed by the
deterministic checkers and statistical models and verified by the author.

\printbibliography

\end{document}
